\documentclass[xcolor={svgnames},
  %hyperref={colorlinks,citecolor=Blue,linkcolor=Blue,urlcolor=DarkBlue}, 
  hyperref={colorlinks}, 
  spanish, 12pt]{beamer}
  \mode<presentation>
  
  \usefonttheme[onlymath]{serif}
  \setbeamertemplate{theorems}[ams style] 
  
  
\usetheme{metropolis}
 \useinnertheme{rectangles} 
\setbeamertemplate{navigation symbols}{}
%\setbeamertemplate{caption}[numbered]
%\useoutertheme{infolines}
\usepackage{cleveref}

\setbeamercovered{highly dynamic}

\newcounter{saveenumi}
\newcommand{\seti}{\setcounter{saveenumi}{\value{enumi}}}
\newcommand{\conti}{\setcounter{enumi}{\value{saveenumi}}}

\resetcounteronoverlays{saveenumi}

%\pagestyle{empty}
\setbeamercolor{emph}{fg=red}
\renewcommand<>{\emph}[1]{%
  {\usebeamercolor[fg]{emph}\only#2{\itshape}#1}%
}

%\usepackage{beamerthemebars}
\usepackage{fontenc}
\usepackage{graphicx}
\usepackage[utf8]{inputenc}
\usepackage[spanish,mexico]{babel}
\usepackage{fontenc}
\usepackage{amsmath}
\usepackage{amsthm}
\usepackage{amssymb}
\usepackage{graphicx}
\usepackage{mathrsfs}
\usepackage{yfonts}
%\usepackage{hyperref}
\usepackage{enumerate}
\usepackage{mathtools}
\usepackage{textcomp}
\usepackage{lmodern}
\usepackage{fancyvrb}
\usepackage{multicol}
\usepackage{color}
\usepackage{verbatim}

\DefineVerbatimEnvironment{ColorVerbatim}{Verbatim}%
  {formatcom=\color{purple},commandchars=\\\{\}}
  
\usepackage{etoolbox}

\BeforeBeginEnvironment{Verbatim}{\begingroup\color{purple}}%
\AfterEndEnvironment{Verbatim}{\endgroup}%
%\usepackage{etoolbox}
% \AtBeginEnvironment{enumerate}{\begin{multicols}{2}}    %%% this line
% \AtEndEnvironment{enumerate}{\end{multicols}}            %%% and this one
% \usepackage{multicol}
%\usepackage[usenames,dvipsnames,svgnames,table]{xcolor}
%\usepackage[urlcolor=blue]{hyperref}
%\numberwithin{section}{part}
\numberwithin{equation}{section} %% Comment out for sequentially-numbered
\numberwithin{figure}{section} %% Comment out for sequentially-numbered

% Automatically generate section title slides in beamer?
% Add support for \subsubsectionpage
\def\subsubsectionname{\translate{}}
\def\insertsubsubsectionnumber{\arabic{subsubsection}}
\setbeamertemplate{subsubsection page}
{
  \begin{centering}
    {\usebeamerfont{subsubsection name}\usebeamercolor[fg]{subsubsection name}\subsubsectionname}%~\insertsubsubsectionnumber}
    \vskip1em\par
    \begin{beamercolorbox}[sep=4pt,center]{part title}
      \usebeamerfont{subsubsection title}\insertsubsubsection\par
    \end{beamercolorbox}
  \end{centering}
}
\def\subsubsectionpage{\usebeamertemplate*{subsubsection page}}

\AtBeginSection{\frame{\sectionpage}}
\AtBeginSubsection{\frame{\subsectionpage}}
\AtBeginSubsubsection{\frame{\subsubsectionpage}}

\theoremstyle{plain}
  \newtheorem{thm}{Teorema}[section]
  \newtheorem{prop}{Proposici\'on}[section]
  \newtheorem{lem}[thm]{Lema}
  \newtheorem{cor}[thm]{Corolario}
  %\newtheorem{rem}{Observaci\'on}[chapter]
  \newtheorem*{sol}{Soluci\'on}
  \newtheorem{alg}{Algoritmo}[section]
  \newtheorem{solved}{Ejercicio Resuelto}
  \newtheorem{evc}{Evaluaci\'on Continua}

\theoremstyle{definition}
  \newtheorem{defn}{Definici\'on}[section]
  \newtheorem{conj}{Conjectura}[section]
  \newtheorem{exmp}{Ejemplo}[section]
  \newtheorem{exe}{Evaluaci\'on Continua}[section]
  \newtheorem{prob}{Problema}[section]
  \newtheorem{rem}{Observaci\'on}[section]
  \newtheorem*{ax}{Axioma}
  \newtheorem{tdv}{Tabla de Verdad}

\theoremstyle{remark}
  \newtheorem{claim}{Afirmaci\'on}[section]
  %\newtheorem{rem}{Observaci\'on}[chapter]
  %\newtheorem*{note}{Nota}
  \newtheorem{case}{Caso}
  \newtheorem{hint}{Sugerencia}[section]


\numberwithin{equation}{section}
\newcommand{\p}{\partial}
%\newcommand{\T}{\mathbb{T}}
\newcommand{\R}{\mathbb{R}}
\newcommand{\N}{\mathbb{N}}
\newcommand{\Q}{\mathbb{Q}}
\newcommand{\abs}[1]{\left|#1\right|}
\newcommand{\f}{\phi}
\newcommand{\vphi}{\varphi}
\newcommand{\vf}{\varphi}
\newcommand{\flow}[2]{\varphi^{#1}\left( #2 \right)}
\renewcommand{\d}[1]{\dot{#1}}
\renewcommand{\r}{\rho}
\newcommand{\A}{\mathcal{A}}
\newcommand{\gam}{\gamma}
\newcommand{\lam}{\lambda}
\renewcommand{\a}{\alpha}
\renewcommand{\b}{\beta}
\newcommand{\om}{\omega}
\newcommand{\iso}{\simeq}
\newcommand{\tensor}{\otimes}
\newcommand{\Z}{\mathbb{Z}}
\newcommand{\set}[1]{\left\{ #1 \right\}}
\newcommand{\inc}{\hookrightarrow}
\renewcommand{\L}{\mathcal{L}}
\renewcommand{\H}{\mathcal{H}}
\newcommand{\D}{\mathcal{D}}
\newcommand{\converge}[1]{\xrightarrow{#1}}
\renewcommand{\cot}{T^{*}}
\newcommand{\cott}[1]{T^{*}\T^{#1}}
\newcommand{\ep}{\epsilon}
\newcommand{\inp}[1]{\langle #1 \rangle}
%\newcommand{\G}{\mathcal{G}}
\newcommand{\hu}{\textbf{h}}
\newcommand{\deck}[1]{\operatorname{Dake}\left( #1 \right)}
\newcommand{\til}[1]{\tilde{#1}}
\newcommand{\Gam}{\Gamma}
\newcommand{\grad}{\nabla}
\newcommand{\del}{\delta}
\newcommand{\id}{\operatorname{Id}}
\newcommand{\Del}{\Delta}
\newcommand{\var}{\Delta}
\newcommand{\avch}[2]{\frac{\Delta #1}{\Delta #2}}
\newcommand{\Avch}[2]{\dfrac{\Delta #1}{\Delta #2}}
\newcommand{\Err}{\operatorname{Err}}
\newcommand{\imply}{\rightarrow}
\newcommand{\wed}{\wedge}
\newcommand{\biconditional}{\longleftrightarrow}
\newcommand{\yields}{\vdash}
\newcommand{\onlyif}{\Rightarrow}
\newcommand{\uset}{\mathbb{U}}
\newcommand{\minus}{\backslash}
\newcommand{\symdif}{\oplus}
\newcommand{\rel}[1]{\boxed{{\color{blue}\textbf{#1}}}}
\newcommand{\dom}[1]{\operatorname{Dom}\left( #1 \right)}
\newcommand{\im}[1]{\operatorname{Im}\left( #1 \right)}
\newcommand{\nrel}[1]{\boxed{{\color{red}\not}{\color{orange}\textbf{#1}}}}
\newcommand{\comp}[2]{#1 \circ #2}

%\date{\today}
\title{Matem\'aticas Discretas \\
Inducci\'on y Recursi\'on}
\author[Juliho Castillo]{\href{https://www.youtube.com/channel/UCb1i-EtybaWWX5urFfmMUWQ}{M. en C. Juliho Castillo}}

\institute[ITESM-CCM]{Tec de Monterrey, Campus Ciudad de M\'exico}
\date{\today}

\begin{document}

\logo{
 \includegraphics[width=1cm,keepaspectratio=true]{./logo.png}
 %LOGO-CECYTEO.png: 200x200 pixel, 72dpi, 7.06x7.06 cm, bb=0 0 200 200
}

\frame{
\titlepage
}

\begin{frame}[allowframebreaks=0.5]
 \tableofcontents
\end{frame}

%\frame{\tableofcontents}

% \AtBeginSubsection[]
% {
%   \begin{frame}
%     \tableofcontents[currentsubsection]
%   \end{frame}
% }

\section{Inducci\'on Matem\'atica}

\subsection{Introducci\'on}

\begin{frame}
 Una propiedad esencial de los naturales $\N=\set{1,2,3,...}$  es la siguiente
 \pause
 
 \begin{ax}[Principio de Inducci\'on Matem\'atica, versi\'on I]
 Sea $P$ una proposici\'on definida en $\N,$ es decir, $P(n)$ toma valores de cierto o falso para cada $n\in \N.$
 
 Supongamos que
 \begin{enumerate}
  \item $P(1)$ es cierto;
  \item $\forall k \in \N: P(k) \onlyif P(k+1).$
 \end{enumerate}

 Entonces $P$ es {cierto} para todo entero positivo $n\in \N.$
 \end{ax}

\end{frame}

\begin{frame}[t]
\begin{exmp}
 Sea $P(n):1+3+5+...+(2n-1)=n^2.$ Demostrar que $P(n)$ es cierta para toda $n \in \N.$ 
\end{exmp}
\end{frame}

\begin{frame}
 \begin{ax}[Principio de Inducci\'on Matem\'atica, versi\'on II]
  Sea $P$ una proposici\'on definida en $\N$ tal que :
  \begin{enumerate}
   \item $P(1)$ es cierta;
   \item $P(k)$ es cierta siempre que $P(j)$ para toda $1\leq j < k.$
  \end{enumerate}
Entonces $P(n)$ es cierta para toda $n\in \N.$
 \end{ax}

\end{frame}

\begin{frame}
 \begin{rem}
  Algunas veces, uno desea demostrar que una proposici\'on es cierta para alg\'un conjunto de enteros
 $$
 \set{a,a+1, a+2,...}
 $$
 donde $a$ es un entero positivo, posiblemente cero. Esto puede hacerse simplemente reemplazando $1$ por $a$ en cualquier versi\'on del Principio de Inducci\'on Matem\'atica. 
 \end{rem}

\end{frame}

\subsection{Notaci\'on ``Sigma''}

\begin{frame}
 La letra griega $\Sigma$ denota adici\'on repetida:
 
 $$
 \sum_{i=a}^{b} f(i)=f(a)+f(a+1)+...+f(b),
 $$ siempre que $a\leq b.$
\end{frame}

\begin{frame}
 \begin{exmp}
  \label{ayr:exmp23.1}
  \begin{enumerate}
   \item $\sum_{j=1}^{5} j = 1+2+3+4+5 =15$
   \item $\sum_{i=0}^{3} \left( 2i+1 \right)=
   1+3+5+7$
   \item $\sum_{i=2}^{10} i^{2}=2^{2}+3^{2}+...+10^{2}$
   \item $\sum_{j=1}^{4}\cos(j\pi)=
   \cos\pi+ \cos 2\pi + \cos 3\pi +\cos 4\pi.$
  \end{enumerate}

 \end{exmp}

\end{frame}

\begin{frame}
\frametitle{Linealidad}
 \begin{prop}
  \label{suma:linealidad}
  \begin{align}
   \sum_{i=a}^{b} cf(i)&=c \sum_{i=a}^{b} f(i)\\
   \sum_{i=a}^{b} f(i)+g(i)&= \sum_{i=a}^{b} f(i)
   +\sum_{i=a}^{b} g(i)
  \end{align}

 \end{prop}

\end{frame}

\begin{frame}
\frametitle{Propiedades}
\begin{align}
 \sum_{k=a}^{b}f(k)&=\sum_{j=a}^{b}f(j)\\
 \sum_{j=a}^{a}f(j)&=f(a) \\
 \sum_{j=a}^{c}f(j)&=\sum_{j=a}^{b}f(j)+\sum_{j=b}^{c}f(j) \\
 \sum_{j=a}^{b+1}f(j)&=\sum_{j=a}^{b}f(j)+f(b)
\end{align}

\end{frame}

\begin{frame}
\begin{exmp}
 Si $f(n)=(2n-1),$ entonces
 $$
 \sum_{i=1}^{n}f(j)=1+3+...+\left( 2n-1 \right)
 $$ es la suma hasta el $n-$\'esimo natural impar. \pause Observe que 
 \begin{enumerate}
  \item $\sum_{j=1}^{1}f(j)=2(1)-1=1.$
  \item $\sum_{i=1}^{n+1}f(j)=\left( \sum_{i=1}^{n}f(j) \right)+\left( 2n+1 \right)$
 \end{enumerate}

\end{exmp}

\end{frame}

\begin{frame}
\begin{exmp}
 Si $f(n)=2^{n-1},$ entonces
 $$
 \sum_{i=1}^{n}f(j)=1+2+...+2^{n-1}
 $$ es la suma de las primeras $n$ potencias de 2 (incluyendo $1=2^{0}$). \pause Observe que 
 \begin{enumerate}
  \item $\sum_{i=1}^{n+1}f(j)=1+2+...+2^{n}$
  \item $\sum_{j=1}^{1}f(j)=2^{1-1}=1.$
  \item $\sum_{i=1}^{n+1}f(j)=\left( \sum_{i=1}^{n}f(j) \right)+2^{n}.$
 \end{enumerate}

\end{exmp}

\end{frame}

\subsection{Ejercicios Resueltos}

\begin{frame}[t]
 \begin{solved}
  Demostrar que $$P(n): 1+2+3+...+n=\frac{1}{2}n\left( n+1 \right)$$
  es cierto para todo $n \in \N.$
 \end{solved}

\end{frame}

\begin{frame}[t]
 \begin{solved}
  Demostrar que $$P(n): 1+2+2^{2}+...+2^{n}=2^{n+1}-1$$
  es cierto para todo $n \in \N.$
 \end{solved}

\end{frame}

\section{Funciones definidas de manera recursiva}

\begin{frame}
Decimos que una funci\'on est\'a \emph{definida recursivamente} si la definici\'on de la funci\'on se refiere a s\'i misma.
\end{frame}

\begin{frame}
Para que la funci\'on est\'e bien definida, debe tener las siguientes dos propiedades:
\begin{enumerate}
 \item Deben existir ciertos argumentos, llamados \emph{valores base,} para los cuales la funci\'on no se refiera a s\'i misma.
 \item Cada vez que la funci\'on se refiera a s\'i misma, el argumento de la funci\'on debe est\'ar m\'as cercano a un valor base.
\end{enumerate}

\end{frame}

\subsection{La funci\'on factorial}

\begin{frame}
El producto de enteros positivos de $1$ hasta $n$ (inclu\'ido) es llamado \emph{$n$ factorial, $n!$}\pause

Es decir, 
$$
n!=n(n-1)\cdots 3\cdot 2 \cdot 1.$$

\end{frame}

\begin{frame}
Por razones combinatorias, es conveniente definir \emph{$0!=1,$} y de esta manera la funci\'on factorial quedar\'a definida para todos los enteros no negativos.
\end{frame}

\begin{frame}[t]
 \begin{rem}
  \begin{enumerate}
   \item $1!=1\cdot0!$\pause
   \item $2!=2\cdot1!$\pause
   \item $3!=3\cdot2!$\pause
   \item $4!=4\cdot3!$
  \end{enumerate}

 \end{rem}

\end{frame}

\begin{frame}
 Es f\'acil observar que para $n \in \N:$
 $$
 n!=n\cdot (n-1)!
 $$
\end{frame}

\begin{frame}
 \begin{defn}[Funci\'on factorial]
  $$n!=
  \begin{cases}
   1 & n=0 \\
   n\cdot(n-1)! & n>0
  \end{cases}
  $$
 \end{defn}

\end{frame}

\begin{frame}
 \begin{rem}
  \begin{enumerate}
   \item El valor de $n!$ factorial esta dado explicitamente para $n=0,$ de manera que $0$ es el valor base. \pause
   \item El valor de $n!, n>0$ est\'a dado en t\'erminos de $n-1,$ que es m\'as cercano al valor base $0.$ \pause   
  \end{enumerate}

  Por tanto, $n!$ es una funci\'on recursiva bien definida.
 \end{rem}

\end{frame}

\begin{frame}[fragile]
 \frametitle{Implentaci\'on iterativa del \emph{factorial} en \texttt{Python}}
 \begin{verbatim}
def factorial(n):
    result = 1
    for i in range(1, n+1):
        result *= i
    return result
 \end{verbatim} 
\end{frame}

\begin{frame}[fragile]
 \frametitle{Implentaci\'on recursiva del \emph{factorial} en \texttt{Python}}
 \begin{verbatim}
def factorial(n):
    z=1
    if n>1:
        z=n*factorial(n-1)
    return z
 \end{verbatim}
 Para m\'as implementaciones, visite \href{https://rosettacode.org/wiki/Factorial}{rosettacode.org/wiki/Factorial}
\end{frame}


\subsection{Suceci\'on de Fibonacci}

\begin{frame}
La celebre sucesi\'on de Fibonacci (usualmente denotada por $F_{0}, F_{1}, F_{2},...$) es como sigue:
$$
0,0,1,2,3,5,8,13,21,34,55,...
$$
\pause
Es decir, $F_{0}=0$  $F_{1}=1$ y cada t\'ermino sucesor es la suma de los dos precedentes.
\end{frame}

\begin{frame}
Por ejemplo, los siguientes dos t\'erminos de la sucesi\'on son
$$34+55=89 \texttt{ y }55+89=144.$$
\end{frame}

\begin{frame}
\begin{defn}[Sucesi\'on de Fibonacci]
$$
F_{n}=
\begin{cases}
 n & n=0,1 \\
 F_{n}=F_{n-2}+F_{n-1} & n>1
\end{cases}
$$
\end{defn}
\end{frame}

\begin{frame}
Este ejemplo es una funci\'on recursiva bien definida, ya que la funci\'on hace referencia a s\'i misma, cuando se usan $ F_{n-2}$ y $F_{n-1},$ y
\begin{enumerate}
 \item los valores base son $0$ y $1;$
 \item los valores de $F_{n}$ est\'an definidos en t\'erminos de valores m\'as peque\~nos $n-2$ y $n-1$ que son m\'as cercanos a los valores base.
\end{enumerate}

\end{frame}

\begin{frame}[fragile]
 \frametitle{Implentaci\'on iterativa de \emph{Fibonacci} en \texttt{Python}}
 \begin{verbatim}
def fibIter(n):
    if n < 2:
        return n
    fibPrev = 1
    fib = 1
    for num in xrange(2, n):
        fibPrev, fib = fib, fib + fibPrev
    return fib
 \end{verbatim} 
\end{frame}

\begin{frame}[fragile]
 \frametitle{Implentaci\'on recursiva de \emph{Fibonacci} en \texttt{Python}}
 \begin{verbatim}
def fibRec(n):
    if n < 2:
        return n
    else:
        return fibRec(n-1) + fibRec(n-2)
 \end{verbatim}
 Para m\'as implementaciones, visite \href{http://rosettacode.org/wiki/Fibonacci\_sequence}{rosettacode.org/wiki/Fibonacci\_sequence}
\end{frame}

\subsection{La funci\'on de Ackermann}

\begin{frame}
\begin{defn}[Funci\'on (fallida) de Ackermann]
 $$
 A(m,n)=
 \begin{cases}
  n+1 & m=0\\
  A(m-1,n) & m\neq0, n=0 \\
  A(m-1, A(m,n-1)) & m\neq 0, n\neq 0
 \end{cases}
 $$
\end{defn}

\end{frame}

\begin{frame}
\begin{defn}[Funci\'on de Ackermann]
 $$
 A(m,n)=
 \begin{cases}
  n+1 & m=0\\
  A(m-1,1) & m\neq0, n=0 \\
  A(m-1, A(m,n-1)) & m\neq 0, n\neq 0
 \end{cases}
 $$
\end{defn}

\end{frame}

\begin{frame}[fragile]
 \frametitle{Implentaci\'on recursiva de \emph{Ackermann} en \texttt{Python}}
 \begin{verbatim}
def ack2(M, N):
    if M == 0:
        return N + 1
    elif N == 0:
        return ack2(M - 1, 1)
    else:
        return ack2(M - 1, ack2(M, N - 1))
 \end{verbatim}
 Para m\'as implementaciones, visite \href{http://rosettacode.org/wiki/Ackermann\_function}{rosettacode.org/wiki/Ackermann\_function}
\end{frame}

\subsection{Ejercicios Resueltos}

\begin{frame}
\begin{solved}
 Sean $a,b$ enteros positivos, y definamos la siguiente funci\'on de manera recursiva:
$$
Q(a,b)=
\begin{cases}
 0 & a<b \\
 Q(a-b,b)+1 & b \leq a
\end{cases}
$$\pause
\begin{enumerate}
 \item Encuentre (i) $Q(2,5);$ (ii) $Q(12,5)$
 \item ?`Qu\'e es lo que hace esta funci\'on? Encuentre $Q(5861,7)$
\end{enumerate}
\end{solved}
\end{frame}

\begin{frame}
\begin{solved}
 Use la definici\'on de la funci\'on de Ackermann para calcular $A(1,3).$
\end{solved}
\end{frame}

\end{document}